\subsection{window-\texorpdfstring{$6$}{6} 的各向异性矩张量、dyadic 边界全息码与可见/隐藏分裂}\label{subsec:conclusion-window6-moment-ldpc-holography-hidden-split}
在 window-\(6\) 的极小紧 Lie 完成、自旋权超立方体、boundary parity 的 rank-\(3\) 环面二挠解释以及 dyadic 体--界全息的精确屏幕与 LDPC 码结构都已固定之后，还可进一步把循环根云的多重度矩张量、精确边界回放的最小屏幕规模、单轴观测的盲秩以及可见异常模块与隐藏补账自由度之间的维数裂解压缩为同一条窗口定量链。其要点在于：循环根云的 escort 二阶矩并不只在 \(q=1\) 处给出 \(\operatorname{diag}(34,32,28)\)，而是形成一条完全显式的对角张量族；该张量族在 \(q=0\) 处严格各向同性，在 \(q>0\) 时与一阶 escort 磁化共同选出 Cartan 轴，并在有限临界指数处发生唯一一次主轴重排。与此同时，\(64\) 个体相微态也不再只是经 fold 压缩后留下一个 \(21\) 维可见异常扇区，而是在要求精确回放时刚性分裂为 \(21\) 维可见模块与至少 \(43\) 维的隐藏二值补账。

\begin{theorem}[循环根云的各向异性质量向量与质量张量]\label{thm:conclusion-window6-weighted-root-cloud-moment-tensor}
在 \(B_3\) 字典下记
\[
\Phi_6^{\mathrm{cyc}}:=\iota_B(X_6^{\mathrm{cyc}})=\Phi(B_3),
\]
并定义多重度函数
\[
m_6(\alpha):=d_6^{\mathrm{bin}}\bigl(\iota_B^{-1}(\alpha)\bigr)\in\{2,3,4\}.
\]
再定义一阶与二阶矩
\[
\mathfrak m_6:=\sum_{\alpha\in\Phi_6^{\mathrm{cyc}}}m_6(\alpha)\,\alpha,
\qquad
\mathfrak M_6:=\sum_{\alpha\in\Phi_6^{\mathrm{cyc}}}m_6(\alpha)\,\alpha\alpha^{\mathsf T}.
\]
则有精确恒等式
\[
\mathfrak m_6=2e_1,
\qquad
\mathfrak M_6=\operatorname{diag}(34,32,28).
\]
因此，window-\(6\) 的循环根云携带一个规范的各向异性质量张量；其最大特征值满足
\[
\rho(\mathfrak M_6)=\lambda_{\max}(\mathfrak M_6)=34=F_9,
\]
并与 boundary sheet parity 的唯一非零 dyadic 差值常数一致。
\end{theorem}
\begin{proof}
由定理 \ref{thm:conclusion-window6-multiplicity-pinning}，根云按多重度分裂为
\[
\Phi_4:=L_4\cup\{e_1,\pm e_2,\pm e_3\},
\qquad
\Phi_3:=L_3,
\qquad
\Phi_2:=L_2\cup\{-e_1\},
\]
其中
\[
L_2=\{(0,\pm1,\pm1)\},
\qquad
L_3=\{(\pm1,0,\pm1)\},
\qquad
L_4=\{(\pm1,\pm1,0)\}.
\]
于是
\[
\mathfrak m_6
=
4\sum_{\alpha\in\Phi_4}\alpha
+3\sum_{\alpha\in\Phi_3}\alpha
+2\sum_{\alpha\in\Phi_2}\alpha.
\]
由对称性可知 \(\sum_{\alpha\in L_4}\alpha=\sum_{\alpha\in L_3}\alpha=\sum_{\alpha\in L_2}\alpha=0\)，而
\[
\sum_{\alpha\in\{e_1,\pm e_2,\pm e_3\}}\alpha=e_1.
\]
故
\[
\sum_{\alpha\in\Phi_4}\alpha=e_1,
\qquad
\sum_{\alpha\in\Phi_3}\alpha=0,
\qquad
\sum_{\alpha\in\Phi_2}\alpha=-e_1,
\]
从而
\[
\mathfrak m_6=4e_1+3\cdot 0+2(-e_1)=2e_1.
\]

再计算二阶矩。对三组长根层有
\[
\sum_{\alpha\in L_4}\alpha\alpha^{\mathsf T}=\operatorname{diag}(4,4,0),
\qquad
\sum_{\alpha\in L_3}\alpha\alpha^{\mathsf T}=\operatorname{diag}(4,0,4),
\qquad
\sum_{\alpha\in L_2}\alpha\alpha^{\mathsf T}=\operatorname{diag}(0,4,4).
\]
另一方面，
\[
\sum_{\alpha\in\{e_1,\pm e_2,\pm e_3\}}\alpha\alpha^{\mathsf T}
=
\operatorname{diag}(1,2,2),
\qquad
(-e_1)(-e_1)^{\mathsf T}
=
\operatorname{diag}(1,0,0).
\]
故
\[
\sum_{\alpha\in\Phi_4}\alpha\alpha^{\mathsf T}=\operatorname{diag}(5,6,2),
\qquad
\sum_{\alpha\in\Phi_3}\alpha\alpha^{\mathsf T}=\operatorname{diag}(4,0,4),
\qquad
\sum_{\alpha\in\Phi_2}\alpha\alpha^{\mathsf T}=\operatorname{diag}(1,4,4).
\]
于是
\[
\mathfrak M_6
=
4\operatorname{diag}(5,6,2)
+3\operatorname{diag}(4,0,4)
+2\operatorname{diag}(1,4,4)
=
\operatorname{diag}(34,32,28).
\]
最后，由推论 \ref{cor:terminal-foldbin6-boundary-pure-f9-alias}，
boundary 三词的 dyadic 预像对唯一非零差值正是 \(34=F_9\)。因此该常数同时也是 \(\mathfrak M_6\) 的最大特征值。证毕。
\end{proof}

\noindent
上述 \((\mathfrak m_6,\mathfrak M_6)\) 实际上只是 cyclic escort 矩族在 \(q=1\) 处的切片；整个 \(q\)-依赖的二阶张量同样保持完全显式。

\begin{theorem}[window-\texorpdfstring{$6$}{6} 的 cyclic escort 二阶矩张量]\label{thm:conclusion-window6-cyclic-escort-second-moment-tensor}
\leanverified{paper\_conclusion\_window6\_cyclic\_escort\_second\_moment\_tensor}
沿用定理 \ref{thm:conclusion-window6-cyclic-escort-first-moment} 的记号。对任意 \(q\ge 0\)，定义
\[
\mathfrak I_{6,\mathrm{cyc}}(q)
:=
\sum_{w\in X_6^{\mathrm{cyc}}} d_6(w)^q\,\iota_B(w)\iota_B(w)^{\mathsf T}.
\]
则有完全显式的对角公式
\[
\mathfrak I_{6,\mathrm{cyc}}(q)
=
4^q\operatorname{diag}(5,6,2)
+3^q\operatorname{diag}(4,0,4)
+2^q\operatorname{diag}(1,4,4).
\]
等价地，其三个特征值为
\[
\lambda_1(q)=5\cdot 4^q+4\cdot 3^q+2^q,
\]
\[
\lambda_2(q)=6\cdot 4^q+4\cdot 2^q,
\]
\[
\lambda_3(q)=2\cdot 4^q+4\cdot 3^q+4\cdot 2^q.
\]
特别地，
\[
\mathfrak I_{6,\mathrm{cyc}}(0)=10I_3,
\qquad
\mathfrak I_{6,\mathrm{cyc}}(1)=\mathfrak M_6=\operatorname{diag}(34,32,28).
\]
因此，cyclic Lie 云在均匀权重下严格各向同性，而纤维多重度权重一旦开启，便立刻把二阶张量从该各向同性基点推出。
\end{theorem}
\begin{proof}
由定理 \ref{thm:xi-window6-dbin-coordinate-quadruple-stratification}，\(X_6^{\mathrm{cyc}}\) 在 \(B_3\) 字典下按退化度分裂为
\[
\Phi_4=\{e_1,\pm e_2,\pm e_3\}\sqcup\{(\pm1,\pm1,0)\},
\]
\[
\Phi_3=\{(\pm1,0,\pm1)\},
\]
\[
\Phi_2=\{-e_1\}\sqcup\{(0,\pm1,\pm1)\}.
\]
于是逐层计算外积和可得
\[
\sum_{\alpha\in\Phi_4}\alpha\alpha^{\mathsf T}=\operatorname{diag}(5,6,2),
\qquad
\sum_{\alpha\in\Phi_3}\alpha\alpha^{\mathsf T}=\operatorname{diag}(4,0,4),
\]
\[
\sum_{\alpha\in\Phi_2}\alpha\alpha^{\mathsf T}=\operatorname{diag}(1,4,4).
\]
分别乘以对应 escort 权 \(4^q,3^q,2^q\) 后求和，即得
\[
\mathfrak I_{6,\mathrm{cyc}}(q)
=
4^q\operatorname{diag}(5,6,2)
+3^q\operatorname{diag}(4,0,4)
+2^q\operatorname{diag}(1,4,4).
\]
对角矩阵的特征值就是三条对角元，从而得到 \(\lambda_1,\lambda_2,\lambda_3\) 的公式。

令 \(q=0\) 时，
\[
\mathfrak I_{6,\mathrm{cyc}}(0)
=
\operatorname{diag}(5+4+1,\ 6+0+4,\ 2+4+4)
=10I_3.
\]
令 \(q=1\) 时，则
\[
\mathfrak I_{6,\mathrm{cyc}}(1)
=
\operatorname{diag}(34,32,28),
\]
与定理 \ref{thm:conclusion-window6-weighted-root-cloud-moment-tensor} 完全一致。证毕。
\end{proof}

\begin{corollary}[window-\texorpdfstring{$6$}{6} escort 二阶主轴的唯一重排]\label{cor:conclusion-window6-cyclic-escort-axis-reordering}
存在唯一实数 \(q_{6,\mathrm{ax}}\in(0,\infty)\)，使得
\[
\lambda_1(q_{6,\mathrm{ax}})=\lambda_2(q_{6,\mathrm{ax}}),
\]
等价地，
\[
4\cdot 3^{q_{6,\mathrm{ax}}}=4^{q_{6,\mathrm{ax}}}+3\cdot 2^{q_{6,\mathrm{ax}}}.
\]
并且对任意 \(q>0\)，有
\[
0<q<q_{6,\mathrm{ax}}
\Longrightarrow
\lambda_1(q)>\lambda_2(q)>\lambda_3(q),
\]
\[
q>q_{6,\mathrm{ax}}
\Longrightarrow
\lambda_2(q)>\lambda_1(q)>\lambda_3(q).
\]
因此，最小主轴对所有 \(q>0\) 都恒为 \(e_3\)，而主轴则仅在 \(e_1\) 与 \(e_2\) 之间发生一次交换。另一方面，定理 \ref{thm:conclusion-window6-cyclic-escort-first-moment} 表明归一化一阶 escort 磁化对所有 \(q>0\) 都固定在 \(e_1\)；故当 \(q>q_{6,\mathrm{ax}}\) 时，向量序参与二阶主轴严格分离于 \(e_1\) 与 \(e_2\)。
\end{corollary}
\begin{proof}
由定理 \ref{thm:conclusion-window6-cyclic-escort-second-moment-tensor}，
\[
\lambda_1(q)-\lambda_3(q)=3(4^q-2^q)>0
\qquad(q>0),
\]
故最小特征值恒为 \(\lambda_3(q)\)。

再记
\[
g(q):=\lambda_1(q)-\lambda_2(q)=4\cdot 3^q-4^q-3\cdot 2^q.
\]
令
\[
x:=\left(\frac32\right)^q>1,
\qquad
a:=\frac{\log 2}{\log(3/2)}>1.
\]
则 \(g(q)=0\) 等价于
\[
h(x):=4x-x^a-3=0.
\]
由于
\[
h''(x)=-a(a-1)x^{a-2}<0
\qquad(x>0),
\]
函数 \(h\) 在 \((0,\infty)\) 上严格凹。又有
\[
h(1)=0,
\qquad
h'(1)=4-a>0,
\qquad
\lim_{x\to\infty}h(x)=-\infty.
\]
因此除 \(x=1\) 之外，方程 \(h(x)=0\) 在 \(x>1\) 上恰有唯一解；相应地，\(g(q)=0\) 在 \(q>0\) 上恰有唯一解，记为 \(q_{6,\mathrm{ax}}\)。其数值审计为
% AUTO-GENERATED by scripts/exp_window6_cyclic_escort_axis_audit.py
\begin{equation}\label{eq:window6_cyclic_escort_axis_audit}
\begin{aligned}
q_{6,\mathrm{ax}} &\approx 4.338068831778077,\\
4\cdot 3^{q_{6,\mathrm{ax}}}-4^{q_{6,\mathrm{ax}}}-3\cdot 2^{q_{6,\mathrm{ax}}} &\approx 4.2632564145606011e-14,\\
\overline{\mathfrak m}_{6,\mathrm{cyc}}(q_{6,\mathrm{ax}})\cdot e_1 &\approx 0.091438769280761.
\end{aligned}
\end{equation}


最后，由 \(g(q)\) 在该唯一根两侧的符号变化，立即得到 \(\lambda_1\) 与 \(\lambda_2\) 的排序结论。证毕。
\end{proof}

\begin{corollary}[矩张量数据对有序 Cartan 轴系的矩阵化钉扎]\label{cor:conclusion-window6-moment-tensor-axis-pinning}
定义联合稳定子
\[
\operatorname{Stab}_{O(3)}(\mathfrak M_6,\mathfrak m_6)
:=
\bigl\{
g\in O(3):
\ g\mathfrak M_6 g^{\mathsf T}=\mathfrak M_6,\ 
g\mathfrak m_6=\mathfrak m_6
\bigr\}.
\]
则有
\[
\operatorname{Stab}_{O(3)}(\mathfrak M_6,\mathfrak m_6)\cong (\ZZ/2\ZZ)^2.
\]
换言之，仅凭循环根云的加权矩阵数据，便可在不显式调用 boundary 词标签的前提下恢复有序 Cartan 轴系；其中第一轴由谱简单性与 \(\mathfrak m_6=2e_1\) 共同定向，而第二、第三轴仅保留独立变号自由度。
\end{corollary}
\begin{proof}
由定理 \ref{thm:conclusion-window6-cyclic-escort-second-moment-tensor} 在 \(q=1\) 处与定理 \ref{thm:conclusion-window6-cyclic-escort-first-moment} 在 \(q=1\) 处，
\[
\mathfrak M_6=\operatorname{diag}(34,32,28),
\qquad
\mathfrak m_6=2e_1.
\]
由于 \(34>32>28\) 为单谱，任意满足 \(g\mathfrak M_6 g^{\mathsf T}=\mathfrak M_6\) 的正交变换 \(g\) 都必须逐一保持三条特征轴
\[
\RR e_1,\qquad \RR e_2,\qquad \RR e_3.
\]
再由 \(g\mathfrak m_6=\mathfrak m_6\) 与 \(\mathfrak m_6=2e_1\)，可知第一轴不仅被保持，而且方向被固定。故仅剩第二、第三坐标上的独立变号，自然给出
\[
\operatorname{Stab}_{O(3)}(\mathfrak M_6,\mathfrak m_6)\cong (\ZZ/2\ZZ)^2.
\]
证毕。
\end{proof}

\begin{conjecture}[escort--Weyl 重整化原理]\label{conj:conclusion-window6-escort-weyl-renormalization}
设 \(m\) 为偶数，并假定其可见 cyclic 扇区允许某个可审计的 pinned Cartan 嵌入
\[
\alpha_m:X_m^{\mathrm{cyc}}\to \mathfrak h_m^\ast.
\]
记
\[
\mathfrak m_m(q)
:=
\sum_{w\in X_m^{\mathrm{cyc}}} d_m^{\mathrm{bin}}(w)^q\,\alpha_m(w),
\qquad
\mathfrak I_m(q)
:=
\sum_{w\in X_m^{\mathrm{cyc}}} d_m^{\mathrm{bin}}(w)^q\,\alpha_m(w)\alpha_m(w)^{\mathsf T}.
\]
则预期以下四项同时成立：
\begin{enumerate}
  \item \(\mathfrak I_m(q)\) 在 multiplicity pinning 所恢复的 Cartan 标架中对角化；
  \item 其特征值只在有限多个 \(q\)-值上发生交叉，且每个交叉点都由纤维尺寸集合 \(\{d_m^{\mathrm{bin}}(w)\}\) 所生成的指数多项式方程给出；
  \item \(\mathfrak m_m(q)\) 的支撑方向始终落在由 multiplicity defect 选出的有限条轴上；
  \item 对 \(m=6\)，定理 \ref{thm:conclusion-window6-cyclic-escort-first-moment}、定理 \ref{thm:conclusion-window6-cyclic-escort-second-moment-tensor} 与推论 \ref{cor:conclusion-window6-cyclic-escort-axis-reordering} 已给出最小非平凡实例：\(q=0\) 时严格各向同性，\(q>0\) 时出现非零轴向 escort 磁化，二阶主轴仅发生一次重排，并且向量序参与张量序参量可落在不同轴上。
\end{enumerate}
若该原理成立，则有限窗的 multiplicity pinning 与 escort 温度并非彼此平行的后处理，而是在纯有限系统内部共同定义一类可见 Cartan 轴的有限温重整化机制。
\end{conjecture}

\begin{theorem}[window-\texorpdfstring{$6$}{6} 的 dyadic 边界像给出精确 \texorpdfstring{$[576,64,12]$}{[576,64,12]} LDPC 码]\label{thm:conclusion-window6-boundary-image-ldpc-576-64-12}
\leanverified{paper\_conclusion\_window6\_boundary\_image\_ldpc\_576\_64\_12}
令 \(\mathcal Q_1\) 为 \([0,1]^6\) 的一层 dyadic 立方剖分，并把其 \(64\) 个顶维 \(6\)-胞腔规范地标记为
\[
\Omega_6=\{0,1\}^6.
\]
定义模 \(2\) 边界像
\[
C_{\partial,6}
:=
\operatorname{im}\!\bigl(
\partial_6:
C_6(\mathcal Q_1;\FF_2)\to C_5(\mathcal Q_1;\FF_2)
\bigr)
\subset \FF_2^{576}.
\]
则 \(C_{\partial,6}\) 是一个参数精确为
\[
[576,64,12]
\]
的规则 LDPC 码。更具体地，
\[
\dim_{\FF_2}C_{\partial,6}=64,
\qquad
d_{\min}(C_{\partial,6})=12,
\qquad
\operatorname{rate}(C_{\partial,6})=\frac{64}{576}=\frac19,
\]
并且其校验度与变量度分别为
\[
4,
\qquad
10.
\]
\end{theorem}
\begin{proof}
由定理 \ref{thm:spg-dyadic-boundary-image-ldpc} 在 \((m,n)=(1,6)\) 的专门化，
\[
N_{1,6}
=
6(2^1+1)2^{1\cdot(6-1)}
=
6\cdot 3\cdot 32
=
576,
\]
并且
\[
\dim_{\FF_2}C_{\partial,6}=2^{1\cdot 6}=64,
\qquad
\operatorname{rate}(C_{\partial,6})=\frac{2}{6(2+1)}=\frac19.
\]
再由同一定理，校验度恒为 \(4\)，变量度为 \(2(n-1)\)，故在 \(n=6\) 时变量度等于 \(10\)。

最后，由推论 \ref{cor:spg-dyadic-boundary-image-min-distance}，
\[
d_{\min}(C_{\partial,6})=2n=12.
\]
综上即得所述精确参数。证毕。
\end{proof}

\begin{corollary}[体相非零精确边界像的最小支撑为 \texorpdfstring{$12$}{12}]\label{cor:conclusion-window6-boundary-ldpc-minsupport-12}
\leanverified{paper\_conclusion\_window6\_boundary\_ldpc\_minsupport\_12}
在定理 \ref{thm:conclusion-window6-boundary-image-ldpc-576-64-12} 的记号下，任意非零体相顶维链 \(x\in C_6(\mathcal Q_1;\FF_2)\) 的精确边界像 \(\partial_6 x\) 都满足
\[
\mathrm{wt}(\partial_6 x)\ge 12.
\]
等价地，不存在支撑在少于 \(12\) 个 \(5\)-胞腔坐标上的非零精确体相边界像。
\end{corollary}
\begin{proof}
这是推论 \ref{cor:spg-dyadic-boundary-image-min-distance} 在 \((m,n)=(1,6)\) 的直接翻译。证毕。
\end{proof}

\begin{theorem}[极小精确屏幕的规模恰为 \texorpdfstring{$64$}{64}]\label{thm:conclusion-window6-exact-screen-minimal-size-64}
\leanverified{paper\_conclusion\_window6\_exact\_screen\_minimal\_size\_64}
记
\[
\overrightarrow{\mathcal F_1^5}
\]
为 \(\mathcal Q_1\) 的全部取向 \(5\)-胞腔所成的标准基。对任意屏幕
\[
S\subset \overrightarrow{\mathcal F_1^5},
\]
定义部分边界观测算子
\[
f_S:=\Pi_S\circ \partial_6:
C_6(\mathcal Q_1;\ZZ)\longrightarrow \ZZ^S.
\]
若 \(f_S\) 为单射，即 \(S\) 能对全部 \(64\) 个体相顶维胞腔作精确回放，则必有
\[
|S|\ge 64.
\]
并且达到该极小值当且仅当 \(S\) 是相应边界行拟阵的一组基；因此极小精确屏幕的基数恰为
\[
|S|_{\min}=64.
\]
\end{theorem}
\begin{proof}
由定理 \ref{thm:xi-time-part9fa-exact-screen-minimal-size-matroid-basis} 在 \((m,n)=(1,6)\) 的专门化，全部满足 \(\ker f_S=0\) 的屏幕中最小可能基数恰为
\[
2^{mn}=2^6=64,
\]
而达到该极小值的全部屏幕恰构成边界行拟阵的基族。证毕。
\end{proof}

\begin{corollary}[单轴内部屏幕的盲秩恰为 \texorpdfstring{$32$}{32}]\label{cor:conclusion-window6-single-axis-screen-blind-rank-32}
固定任一坐标方向 \(i\in\{1,\dots,6\}\)，并令
\[
S_{\{i\}}^{\mathrm{int}}\subset \overrightarrow{\mathcal F_1^5}
\]
为全部法向平行于 \(e_i\) 的内部 \(5\)-胞腔所组成的单轴内部屏幕。则其最小审计补全成本满足
\[
\min_{(R,r_{S_{\{i\}}^{\mathrm{int}}})}\cdim(R)=32.
\]
因此单轴内部屏幕只能捕获一半的体相线性信息，而剩余另一半必须通过外置审计寄存器补齐。
\end{corollary}
\begin{proof}
由定理 \ref{thm:spg-internal-coordinate-bundle-screen-cost-closed-form}，对于 \(|J|=s\) 的内部坐标束屏幕有
\[
\min_{(R,r_{S_J^{\mathrm{int}}})}\cdim(R)=2^{m(n-s)}.
\]
代入 \(m=1\)、\(n=6\)、\(s=1\) 即得
\[
\min_{(R,r_{S_{\{i\}}^{\mathrm{int}}})}\cdim(R)=2^{1\cdot(6-1)}=32.
\]
证毕。
\end{proof}

\begin{theorem}[window-\texorpdfstring{$6$}{6} 单轴屏幕的半盲刚性]\label{thm:conclusion-window6-oneaxis-screen-halfblind-rigidity}
\leanverified{paper\_conclusion\_window6\_oneaxis\_screen\_halfblind\_rigidity}
在 \(Q_1^6\) 的 window-\(6\) 口径下，设 \(S^{(1)}\) 为任一单轴内部屏幕，并定义其最小精确化代价为
\[
\operatorname{Cost}(S^{(1)}):=\min_{(R,r_{S^{(1)}})}\cdim(R).
\]
则有
\[
\operatorname{Cost}(S^{(1)})=2^{1\cdot(6-1)}=32.
\]
另一方面，全部极小精确屏幕的规模恰为
\[
|S|_{\min}=2^6=64.
\]
因此，单轴内部屏幕恰丢失了整整一半的 bulk 自由度；任何单轴方案若不额外补入 \(32\) 个独立边界坐标，便不可能达到精确回放。
\end{theorem}
\begin{proof}
由推论 \ref{cor:conclusion-window6-single-axis-screen-blind-rank-32}，任一单轴内部屏幕的最小审计补全成本满足
\[
\operatorname{Cost}(S^{(1)})=32.
\]
另一方面，定理 \ref{thm:conclusion-window6-exact-screen-minimal-size-64} 已证明：若屏幕能对全部 \(64\) 个体相顶维胞腔作精确回放，则其规模至少为 \(64\)，且该下界可达。两式合并即得结论。证毕。
\end{proof}

\begin{corollary}[window-\texorpdfstring{$6$}{6} 一轴精确回放的二尺度障碍]\label{cor:conclusion-window6-oneaxis-exact-replay-two-scale-obstruction}
沿用定理 \ref{thm:conclusion-window6-oneaxis-screen-halfblind-rigidity} 的记号。
任意非零体相顶维链 \(x\in C_6(\mathcal Q_1;\FF_2)\) 的精确边界像都满足
\[
\mathrm{wt}(\partial_6 x)\ge 12.
\]
因此，window-\(6\) 的一轴精确回放同时受两条彼此独立的硬障碍支配：
\[
\operatorname{Cost}(S^{(1)})=32,
\qquad
\min\{\mathrm{wt}(\partial_6 x):x\neq 0\}=12.
\]
换言之，单轴失明并非少量局部校验即可修补的缺陷；它同时要求 \(32\) 个独立补账方向，并且每一个非平凡精确修正都不可压缩到少于 \(12\) 个边界坐标。
\end{corollary}
\begin{proof}
第一条即定理 \ref{thm:conclusion-window6-oneaxis-screen-halfblind-rigidity}。
第二条由推论 \ref{cor:conclusion-window6-boundary-ldpc-minsupport-12} 直接给出。
两者量纲不同，分别控制全局自由度亏损与局部修正最小支撑，因而彼此不能替代。证毕。
\end{proof}


\leanverified{paper\_conclusion\_window6\_visible\_hidden\_binary\_splitting\_21\_43}
\begin{proposition}[可见异常模块之外至少需要 \texorpdfstring{$43$}{43} 维隐藏二值补账]\label{prop:conclusion-window6-visible-hidden-binary-splitting-21-43}
记
\[
B_6:=C_6(\mathcal Q_1;\FF_2)\cong \FF_2^{64},
\]
并在 \(\ZZ/2\ZZ=\FF_2\) 的自然识别下，把 fold-gauge 口径下的可见异常模块记为
\[
V_6:=\bigl(\mathcal G_6^{\mathrm{bin}}\bigr)^{\mathrm{ab}}
\cong
\FF_2^{21}.
\]
若存在 \(\FF_2\)-线性精确因子化
\[
B_6\xrightarrow{\ \iota\ }V_6\oplus R
\xrightarrow{\ \pi\ }B_6,
\qquad
\pi\circ\iota=\operatorname{id}_{B_6},
\]
则必有
\[
\dim_{\FF_2}R\ge 43.
\]
因此，在坚持精确体相重建的前提下，\(21\) 维可见异常模块之外至少还需 \(43\) 个隐藏二值自由度。
\end{proposition}
\begin{proof}
由定理 \ref{thm:window6-abelianized-parity-charge-root-cartan-splitting} 或推论 \ref{cor:conclusion-circle-dimension-anomaly-rank-budget-splitting}，可见异常模块的维数恰为 \(21\)。另一方面，由
\[
\pi\circ\iota=\operatorname{id}_{B_6}
\]
可知 \(\iota\) 必为单射，故
\[
64=\dim_{\FF_2}B_6
\le
\dim_{\FF_2}(V_6\oplus R)
=
21+\dim_{\FF_2}R.
\]
于是
\[
\dim_{\FF_2}R\ge 64-21=43.
\]
证毕。
\end{proof}

\begin{theorem}[boundary parity 的最小 torus holonomy 秩恰为 \texorpdfstring{$3$}{3}]\label{thm:conclusion-window6-boundary-parity-minimal-torus-rank}
\leanverified{paper\_conclusion\_window6\_boundary\_parity\_minimal\_torus\_rank}
记 window-\(6\) 的 boundary parity 群为
\[
Z_6^{\mathrm{bd}}\cong (\ZZ/2\ZZ)^3.
\]
定义其最小 torus holonomy 秩
\[
r_{\mathrm{tor}}(Z_6^{\mathrm{bd}})
:=
\min\Bigl\{
r:\ \exists\ \text{忠实嵌入 }Z_6^{\mathrm{bd}}\hookrightarrow \TT^r
\Bigr\}.
\]
则有
\[
r_{\mathrm{tor}}(Z_6^{\mathrm{bd}})=3.
\]
因此，boundary 拓扑扇区既不能压缩为单一相位，也不能压缩为两个独立相位；它要求且仅要求三条独立的紧 holonomy 通道。
\end{theorem}
\begin{proof}
由定理 \ref{thm:conclusion-window6-boundary-torus-2torsion}，
\[
Z_6^{\mathrm{bd}}\cong T_6[2],
\qquad
T_6\cong \TT^3,
\]
故显然有忠实嵌入 \(Z_6^{\mathrm{bd}}\hookrightarrow \TT^3\)，从而
\[
r_{\mathrm{tor}}(Z_6^{\mathrm{bd}})\le 3.
\]

反过来，设 \(Z_6^{\mathrm{bd}}\hookrightarrow \TT^r\) 为任一忠实嵌入。由于 \(Z_6^{\mathrm{bd}}\) 中元素均为二阶元，其像必落在
\[
(\TT^r)[2]\cong (\ZZ/2\ZZ)^r
\]
之中。故 \((\ZZ/2\ZZ)^3\) 作为子群嵌入 \((\ZZ/2\ZZ)^r\) 必要求 \(r\ge 3\)。于是
\[
r_{\mathrm{tor}}(Z_6^{\mathrm{bd}})\ge 3.
\]
合并两侧不等式即得结论。证毕。
\end{proof}

\begin{conjecture}[window-\texorpdfstring{$6$}{6} 的规范 CSS 严格化]\label{conj:conclusion-window6-css-strictification}
为避免与 dyadic 立方复形 \(\mathcal Q_1\) 混淆，记 \(\mathsf Q_6^{\mathrm{CSS}}\) 为一个待构造的 window-\(6\) 规范 CSS 严格化对象。猜想存在这样的 \(\mathsf Q_6^{\mathrm{CSS}}\)，使其同时满足：
\begin{enumerate}
  \item 其 \(Z\)-型经典阴影恰为定理 \ref{thm:conclusion-window6-boundary-image-ldpc-576-64-12} 的边界码
  \[
  C_{\partial,6}\subset \FF_2^{576};
  \]
  \item 其逻辑二值模 \(L_6\) 满足
  \[
  L_6\cong \FF_2^{64},
  \]
  并存在满射
  \[
  L_6\twoheadrightarrow V_6\cong \FF_2^{21}
  \]
  ，其核是一个 \(43\) 维隐藏规范子空间；
  \item 其紧连续 holonomy 的最小 torus 秩等于
  \[
  r_{\mathrm{tor}}(Z_6^{\mathrm{bd}})=3;
  \]
  \item 其低能可见连续极限的各向异性二次型由定理 \ref{thm:conclusion-window6-weighted-root-cloud-moment-tensor} 的质量张量
  \[
  \mathfrak M_6=\operatorname{diag}(34,32,28)
  \]
  控制。
\end{enumerate}
若此猜想成立，则 window-\(6\) 将给出一个完全精确的混合全息对象：原始体相逻辑模为 \(64\) 维，精确边界码为 \([576,64,12]\)，可见异常模为 \(21=18+3\) 维，隐藏规范模为 \(43\) 维，紧连续 holonomy 的最小秩为 \(3\)，而可见连续谱的主惯量则由 \(34=F_9\) 锁定。
\end{conjecture}

\paragraph{统一读法}
上述链条可压缩为
\[
B_6=C_6(\mathcal Q_1;\FF_2)\cong \FF_2^{64}
\longrightarrow
C_{\partial,6}\ \text{为}\ [576,64,12]\ \text{边界 LDPC 码}
\longrightarrow
|S|_{\min}=64
\longrightarrow
\operatorname{Audit}(S_{\{i\}}^{\mathrm{int}})=32
\longrightarrow
64=21+43
\longrightarrow
r_{\mathrm{tor}}(Z_6^{\mathrm{bd}})=3
\longrightarrow
\mathfrak M_6=\operatorname{diag}(34,32,28).
\]
其几何内容是循环根云主轴与 dyadic 屏幕的共同钉扎；其编码内容是精确边界像的 \([576,64,12]\) 结构、最小屏幕 \(64\) 与单轴盲秩 \(32\)；其代数内容则是可见异常模块与隐藏补账自由度的 \(21/43\) 维裂解，以及 boundary parity 的最小 torus 实现秩 \(3\)。而 \(34=F_9\) 的双重出现则把 boundary sheet parity 的离散差值与循环根云质量张量的主特征值压缩为同一个谱常数。

\endinput
